\documentclass[11pt]{article}
\usepackage{nopageno}
\setcounter{secnumdepth}{1}
\usepackage[margin=1in]{geometry}
\usepackage{enumitem}
\begin{document}
\pagestyle{empty}

\noindent \textbf{Synergistic Activities $|$ Dr.\ Rebecca Napolitano}

\vspace{0.3cm}

\begin{enumerate}[leftmargin=*, labelsep=0.5em, itemsep=0.4cm]

\item \textbf{Mentoring and Developing the Next Generation of Researchers:}
Lead the AE Research Scholars undergraduate research program in the
Architectural Engineering department at Penn State, mentoring 5 undergraduates
per year on impactful projects such as analyzing Montpelier flood data.
Mentored 2 undergraduate students through the University of Washington's REU
program from 2022--2023, with projects focused on post-tornado historic masonry
building reconstruction.

\item \textbf{Innovations in Teaching and Curriculum Development:}
Co-leading DOE- and NSF-funded projects on contextualizing data science
education for engineering students, with the potential to transform how these
critical skills are taught. The project involves assessing how problem-based
learning and real-world case studies can make data science more relevant and
engaging for engineering students. Preliminary findings, in review at the
\textit{Journal of Engineering Education} and the \textit{Journal for Civil
Engineering Education}, and presented at the ASEE Annual Conference, suggest
that this contextualized approach significantly improves student learning
outcomes and attitudes towards data science. Recognized through high student
evaluations (4.9/5.0 average over 5 semesters). Invited presenter at the
Graduate Women in Engineering seminar series.

\item \textbf{Development and Dissemination of Open-Source Data and Models:}
Made research data from the 2021 Quad-State Tornado freely available through
the NHERI DesignSafe-CI platform. This has allowed other researchers to build
upon this work and accelerate progress in this area. Made educational materials
from courses and grants freely available through a website. This website
provides information about a data science program offering six courses tailored
for engineering students. These courses cover a range of crucial topics,
including advanced data science, engineering ethics, data management, machine
learning, parametric modeling, and an introduction to data science for
engineers. This comprehensive curriculum equips engineering students with the
essential skills and knowledge required to excel in the data-driven world,
addressing both technical aspects and ethical considerations within the field
of data science.

\item \textbf{Building Interdisciplinary Research Collaborations:}
Established a collaboration with Dr.\ Guangqing Chi of Rural Sociology and
Demography at Penn State on an NSF-funded project to develop community-driven
innovations for resilient infrastructure in rural Alaska Native communities.
Dr.\ Chi provides guidance on community engagement strategies, participatory
research methods, and social science methodologies to ensure the project is
responsive to local needs and contexts. This collaboration has resulted in
several joint publications and a workshop for Alaska Native community grade
6--12 students on infrastructure monitoring and maintenance.

\item \textbf{Service to the Profession and Broader Impacts:}
Vice President of the International Scientific Committee on the Analysis and
Restoration of Structures of Architectural Heritage (ISCARSAH), leading
emerging professionals initiatives and chairing the Climate Adaptation Working
Group to develop best practice guidelines. Editorial board member for the
\textit{Journal of Civil Structural Health Monitoring}. Regularly serve as a
peer reviewer for 10 journals and 3 conferences in the field. Organize
conference sessions and workshops, such as the Structural Analysis of Heritage
Structures Mini Symposium at the Engineering Mechanics Institute conference
(2019--2024).

\end{enumerate}

\end{document}
